\subsection{Messverfahren}
Grundaufbau jedes Versuchs ist eine optische Bank, auf der eine Lichtquelle, verstellbare Linsen
und Blenden montiert sind. Es werden bikonvexe und achromatische Linsen verwendet, um optische Abbildungen zu untersuchen. 
Die vier Aufgaben sind:
\begin{enumerate}
    \item Parameter eines Achromats: um mit der achromatischen Linse scharfe Bilder auf einem Schirm zu erzeugen und damit zu bestimmen. Der Objektabstand g wird innerhalb bestimmter Bereiche variiert, um zu beobachten, wie sich das Bild und der Bildabstand b verändern (scharfgestellte Position des Schirms wird gesucht) 
    \item Brennweitenbestimmung nach Besselverfahren
    \item Untersuchung chromatischer und sphärischer Abberation durch Farbfilter bzw. durch verschiedene Blenden
    \item Bestimmung des Auflösungsvermögens eines einfachen Mikroskops   
\end{enumerate}

\xfigure{htbp}{width=0.7\textwidth}{Aufbau.png}{Versuchsaufbau: optische Bank\autocite{Skript_1}}{optische Bank mit Apparaturen \autocite{Skript_1}}

\subsection{Messprotokoll}
\includegraphics[page=1,width=\textwidth]{V31.pdf}
\newpage
\includegraphics[page=2,width=\textwidth]{V31.pdf}
\newpage
\includegraphics[page=3,width=\textwidth]{V31.pdf}
